% Abhakseite „Das kann ich“ – Zone nur im Gesamt (\mitzone)
%% TODO Ich-kann-Sätze: Platzhalter wie in den Aufgabendateien
\begin{abhakseite}
\ifdefined\mitzone
\abhakgruppe{Kennst du schon}
\abhak{1}{Lineare Gleichungssysteme lösen}
\abhak{2}{Lineare Funktionen}
\abhak{3}{Parabelformen kennen}
\abhak{4}{Symmetrie am Term erkennen}
\abhak{5}{Ableitungsregeln anwenden}
\abhak{6}{Sinus- und Kosinusgraph mit Periode und Amplitude}
\abhak{7}{Lineare Gleichungssysteme lösen – Fehler finden}
\abhak{8}{Lineare Gleichungssysteme lösen – selbst rechnen}
\fi
\abhakgruppe{Einheit 1 · Ansatz und Punktbedingungen}
\abhak{9}{Ansatz und Punktbedingungen}
\abhak{10}{Ansatz und Punktbedingungen – Fehler finden, Begründen}
\abhakgruppe{Einheit 2 · Bedingungen mit Ableitung}
\abhak{11}{Bedingungen mit Ableitung}
\abhak{12}{Zwei Parameter einer Exponentialfunktion aus Funktionswert und Änderungsrate an einer Stelle bestimmen}
\abhak{13}{Bedingungen mit Ableitung – Fehler finden, Begründen}
\abhakgruppe{Einheit 3 · Sonderansätze und Modellkritik}
\abhak{14}{Sonderansätze und Modellkritik}
\abhak{15}{Zwei Parameter einer Funktion aus einer vorgegebenen Bogenlängenformel und einem Punkt bestimmen}
\abhak{16}{Existenz einer quadratischen Funktion zu vier Wert- und Steigungsbedingungen über das überbestimmte Gleichungssystem untersuchen}
\abhak{17}{Unmöglichkeit waagerechter Tangenten an den Rändern über die Ableitung mit Parametern begründen}
\abhak{18}{Sonderansätze und Modellkritik – Fehler finden, Begründen}
\end{abhakseite}
